\chapter{Dominions und der Zigzag-Satz}

\section{Bestimmung durch eine Unterhalbgruppe}

\begin{definition}[Dominion]
Für eine Unterhalbgruppe $U\subseteq S$ ist ihr \emph{Dominion} in $S$
die Menge $\operatorname{Dom}_S(U)$ aller $d\in S$ mit folgender
Eigenschaft: Für jede Halbgruppe $T$ und alle Homomorphismen
$f,g:S\to T$ gilt
\[
 f|_U=g|_U\ \Longrightarrow\ f(d)=g(d).
\]
\end{definition}

\begin{theorem}[Isbells Zigzag-Satz]
\label{b28guide:thm:zigzag}
Für jede Unterhalbgruppe $U\subseteq S$ und jedes $d\in S$ gilt:
$d\in\operatorname{Dom}_S(U)$ genau dann, wenn $d\in U$ ist oder
es eine Zahl $m\geq1$, Elemente $a_0,\ldots,a_{2m}\in U$
und $x_1,\ldots,x_m,y_1,\ldots,y_m\in S$ gibt mit
\[
 \begin{aligned}
 d&=a_0y_1=x_ma_{2m},\\
 a_0&=x_1a_1,& a_{2m-1}y_m&=a_{2m},
 \end{aligned}
\]
und, für $1\leq i<m$,
\[
 \begin{aligned}
 a_{2i-1}y_i&=a_{2i}y_{i+1},\\
 x_i a_{2i}&=x_{i+1}a_{2i+1}.
 \end{aligned}
\]
Diese endliche Gleichungsanordnung heißt \emph{Zigzag}.
Für $m=1$ entfällt die zweite Gleichungsgruppe.
Weder $U$ noch $S$ müssen endlich sein.
\cite{isbell1966,alam2019}
\end{theorem}

\section{Dominions und Amalgamierung}

\begin{theorem}[Dominion als Schnitt zweier Kopien]
\label{b28guide:thm:dominion-amalgam}
Für jede Unterhalbgruppe $U\subseteq S$ seien $S_+,S_-$ disjunkte
Kopien von $S$, mit den entsprechenden Elementen $s_+,s_-$.
Bilde
\[
 P=(S_+*S_-)/\rho,\qquad
 \rho=\operatorname{Cg}_{S_+*S_-}
 \bigl(\{(u_+,u_-):u\in U\}\bigr).
\]
Die kanonischen Homomorphismen $j_+,j_-:S\to P$ sind Einbettungen.
Für alle $s,t\in S$ gilt
\[
 j_+(s)=j_-(t)
 \iff s=t\ \text{und}\ s\in\operatorname{Dom}_S(U).
\]
Insbesondere ist
\[
 j_+(S)\cap j_-(S)
 =j_+\bigl(\operatorname{Dom}_S(U)\bigr)
 =j_-\bigl(\operatorname{Dom}_S(U)\bigr).
\]
Zwei Kopien von $S$ lassen sich daher genau über $U$ zusammenfügen,
ohne zusätzliche Identifizierungen zwischen den Kopien,
genau dann, wenn $\operatorname{Dom}_S(U)=U$ gilt.
Die Injektivität gilt hier für die beiden Kopien derselben Halbgruppe;
allgemeine Pushouts aus Satz~\ref{b28guide:thm:pushout} können sie verlieren.
\cite{isbell1966,howie1968}
\end{theorem}

\section{Kategorische Epimorphismen}

\begin{definition}[Epimorphismus]
Ein Homomorphismus $\varphi:S\to T$ heißt hier \emph{kategorischer
Epimorphismus}, wenn für jede Halbgruppe $V$ und alle Homomorphismen
$f,g:T\to V$ gilt
\[
 f\circ\varphi=g\circ\varphi\ \Longrightarrow\ f=g.
\]
Dies ist die rechtsseitige Kürzbarkeit von $\varphi$ unter Komposition;
es ist keine Kürzbarkeitsannahme an die Multiplikation von $S$ oder $T$.
\end{definition}

\begin{theorem}[Dominions und Epimorphismen]
\label{b28guide:thm:epi-dominion}
Für jede Unterhalbgruppe $U\subseteq S$ ist
$\operatorname{Dom}_S(U)$ eine Unterhalbgruppe, die $U$ enthält.
Es gilt $\operatorname{Dom}_S(\varnothing)=\varnothing$.
Für jeden Homomorphismus $\varphi:S\to T$ gilt
\[
 \varphi\text{ ist ein kategorischer Epimorphismus}
 \iff
 \operatorname{Dom}_T(\varphi(S))=T .
\]
Jeder surjektive Homomorphismus ist ein kategorischer Epimorphismus.
Die Umkehrung gilt nicht allgemein: Es gibt nichtsurjektive
kategorische Epimorphismen von Halbgruppen.
\cite{isbell1966,alam2019}
\end{theorem}
